\section{本章小结}
\label{sec:cdist-summary}

本章把前两章的分布表示和数值特征落实到十二个可以直接计算的模型上。离散模型都由试验机制推出：一次成败给出Bernoulli分布；固定次数的独立重复给出二项分布；有限总体的不放回抽样给出超几何分布；等待首次成功给出几何分布；机会多、每次稀少的计数，作为二项分布的极限，给出Poisson分布。连续模型则由取值范围和形状要求确定，先验证密度的归一化，再计算矩。

几个计算工具贯穿全章。指示变量把计数写成事件之和，期望由线性性直接得到，不需要独立性；方差要用到两两交事件的概率，独立与否在这里起决定作用，不放回抽样的负牵连带来有限总体修正。尾概率求和处理了取值无穷多的离散变量。Gauss积分、Gamma函数与Beta函数提供了连续模型的归一化常数，分部积分与代换给出各阶矩。

每个公式都带着条件。几何分布要求$p>0$，并且要分清计数是否包含成功的那一次；超几何方差的公式要求$N>1$；Gamma分布有率和尺度两种写法，同一组数字按两种写法理解，期望可以相差好几倍；t分布只有阶数小于$\nu$的矩，$\nu\leq1$时期望不存在，Cauchy分布的对称中心是中位数而不是期望。尺度参数与标准差的比例因分布而异，只有高斯分布的$\sigma$恰为标准差。

模型的比较说明，均值相同甚至均值和方差都相同，也不能让一个模型代替另一个：试验机制、支撑、方差随均值的变化规律和尾部形状都可能不同。方差相同的分布，远处的尾概率可以相差几个数量级。下一章研究这些尾概率：已知完整分布时直接计算，只知道均值、方差或更高阶矩时，能给出怎样的上界，本章的模型将用来检验这些上界的松紧。
